\documentclass{beamer}

\usepackage[utf8]{inputenc}
\usepackage[spanish]{babel}

\usetheme{Rochester}
\usecolortheme{seagull}
\beamertemplatenavigationsymbolsempty

\title[Inflación y poder adquisitivo]{Inflación heterogénea\\y poder adquisitivo en Argentina}
\subtitle{Cuando los precios suben, ¿todos pierden por igual?}
\author{Javier Spina}
\institute{ECyT - UNSAM}
\date{27 de marzo de 2026}

\begin{document}

\begin{frame}
    \titlepage
\end{frame}


\section{Introducción}
\begin{frame}{Introducción}
    \begin{itemize}
        \item La inflación en Argentina es \textbf{persistente y recurrente}: un fenómeno que atraviesa el debate público y la vida cotidiana.
        \item El IPC oficial pondera \textbf{``un hogar promedio'' que no existe}.
        \item Quien gasta el 50\% en alimentos enfrenta una inflación
        \textbf{distinta} a quien gasta el 50\% en servicios o alquiler.
    \end{itemize}
\end{frame}


\section{Contexto del problema}
\begin{frame}{Contexto del problema}
    \begin{itemize}
        \item Historia de alta inflación y shocks económicos.
        \item Las estructuras de gasto varían y evolucionan de forma 
        \textbf{distinta} según el perfil de cada trabajador.
        \item Los sectores de menores ingresos pierden poder adquisitivo 
        en ambos contextos: en \textbf{alta inflación} porque sus canastas 
        suben más, en \textbf{baja inflación} porque el IPC no representa 
        su gasto.
        \item La metodología de medición del INDEC se basa en canastas 
        de 2004. La actualización prevista \textbf{no se implementó}.
        \item Muchos asalariados formales estan cayendo bajo la linea de pobreza, incluso midiendo con el IPC oficial.
    \end{itemize}
    \vspace{0.5em}
    \begin{block}{}
        Hay datos para estudiar esto con mayor precisión.
    \end{block}
\end{frame}

\section{Ejemplos y casos de uso}
\begin{frame}{Ejemplos y casos de uso}

    \begin{table}
        \centering
        \small
        \begin{tabular}{ll}
            \hline
            \textbf{Perfil} & \textbf{Gasto dominante}\\
            \hline
            Trabajadora informal, NOA & Alimentos\\
            Empleado formal, Patagonia & Vivienda y servicios\\
            Docente universitario, CABA & Alimentos y alquiler\\
            Profesional independiente, Cuyo & Servicios y esparcimiento\\
            \hline
        \end{tabular}
    \end{table}

    \vspace{0.5em}
    \begin{block}{}
        El mismo IPC oficial para los cuatro. ¿Tiene sentido?
    \end{block}

\end{frame}

\section{Definición del problema}
\begin{frame}{Definición del problema}
    \begin{block}{Pregunta central}
        ¿Cuál es la tasa de inflación real que enfrenta cada perfil de trabajador argentino, y qué tan vulnerable es su situación cuando se mide con ese índice en lugar del IPC oficial?
    \end{block}
    \vspace{0.5em}
    \begin{itemize}
        \item Deflactar salarios con un único IPC distorsiona quién pierde y cuánto.
        \item Cada trabajador tiene su propia canasta, \textbf{su propia inflación}.
        \item Clustering para \textbf{descubrir los perfiles} que emergen de los datos.
    \end{itemize}
\end{frame}


\section{Selección de papers}
\begin{frame}{Selección de papers}
    \begin{itemize}
        \item \textbf{Jaravel (2024)} --- \textit{Distributional Consumer Price Indices}\\
        Construye D-CPIs por perfil sociodemográfico y muestra que la desigualdad real creció un 45\% más rápido de lo que indica el IPC oficial.
        \vspace{0.4em}
        \item \textbf{Lódola, Busso y Cerimedo (2000)} --- \textit{Sesgos en el IPC: el sesgo plutocrático en Argentina}\\
        IPC por decil en alta inflación. La inflación alta tuvo comportamiento anti-rico. La inflación nominal baja tuvo comportamiento anti-trabajador.
        \vspace{0.4em}
        \item \textbf{Deaton (1985)} --- \textit{Panel data from time series of cross-sections}\\
        Fundamento metodológico del seguimiento de cohortes poblacionales y su evolución en el tiempo a traves de la media del grupo (pseudo-panel).
    \end{itemize}
\end{frame}

\end{document}
